\hypertarget{spec-v3-public-release-v3-from-personal-tool-to-public-self-hostable-server-capable-project}{%
\subsection{v3: from personal tool to public, self-hostable,
server-capable
project}\label{spec-v3-public-release-v3-from-personal-tool-to-public-self-hostable-server-capable-project}}

Supersedes v1/v2 wherever they conflict. Keeps: streak algorithm (v1
§8), pack format and levels/milestones (v2 §3, §8b), verifier interface
shape (v2 §7), day-wise/level logic. Changes: the project is no longer
about one person or one subject.

\hypertarget{spec-v3-public-release-0-naming-your-call--pick-one-or-suggest-your-own-update-throughout-before-publishing}{%
\subsubsection{0. Naming (your call --- pick one or suggest your own,
update throughout before
publishing)}\label{spec-v3-public-release-0-naming-your-call--pick-one-or-suggest-your-own-update-throughout-before-publishing}}

Working name used below: \textbf{Ascend}. Binary: \texttt{ascend}.
Alternatives: \textbf{Pathforge}, \textbf{Wayrune}. Rationale: generic,
not tied to hacking, describes climbing levels toward mastery. Rename is
a find-and-replace; nothing below depends on the word.

\hypertarget{spec-v3-public-release-1-what-changes-and-why}{%
\subsubsection{1. What changes, and why}\label{spec-v3-public-release-1-what-changes-and-why}}

\begin{longtable}[]{@{}P{.24\textwidth}P{.35\textwidth}P{.33\textwidth}@{}}
\toprule\noalign{}
& Before (hackstreak, personal) & Now (Ascend, public) \\
\midrule\noalign{}
\endhead
\bottomrule\noalign{}
\endlastfoot
Who it\textquotesingle s for & You & Anyone learning anything, you
included \\
Subject & Baked toward cyber & Domain-agnostic; cyber is one example
pack among many \\
Where it runs & Your laptop only & Standalone on one machine, OR a
server (e.g. your Pi) with many clients \\
Your data & In the same repo as the code & In your own private
packs/profile, outside the installable package \\
Users & Implicitly one & Zero or more, with isolated data, if you run it
as a server \\
Nudge source & A sync hack to the Pi & The server (if you run one) ---
no sync needed, it\textquotesingle s just the server \\
\end{longtable}

Compare to the projects you named: \textbf{Trilium} and
\textbf{CherryTree} both work great as a single local app, and
\textbf{separately} offer (Trilium) or could offer (CherryTree via file
sync) a server/sync mode for multiple devices. Ascend follows the same
shape: the core engine and CLI work with zero setup and zero
dependencies on one machine; running a server is an opt-in upgrade, not
a requirement.

\hypertarget{spec-v3-public-release-2-two-modes-one-engine}{%
\subsubsection{2. Two modes, one engine}\label{spec-v3-public-release-2-two-modes-one-engine}}

\begin{verbatim}
engine/           <- pure logic: streak math, level/gate logic, pack parsing, verifier interface. No I/O assumptions.
storage/          <- SQLite access. Same schema whether local or server.
cli/              <- `ascend` command. Talks to storage/ DIRECTLY in standalone mode,
                      or to a remote server over HTTP in client mode (one `--server URL` / config line switches it).
server/           <- optional. A thin HTTP API in front of engine/ + storage/. Adds multi-user + nudging + a web view.
\end{verbatim}

\begin{itemize}
\tightlist
\item
  \textbf{Standalone mode} (default, what v1/v2 described):
  \path|pip install ascend-tracker|, stdlib only, one local SQLite
  file. Nothing to deploy. This stays the easiest path and the one the
  README leads with.
\item
  \textbf{Server mode} (new, opt-in): run \texttt{ascend-server} (adds
  \texttt{fastapi}+\texttt{uvicorn}, server-side only --- the CLI client
  stays dependency-free) somewhere always-on --- your Pi, a VPS, a spare
  machine. Point any number of CLI installs or a browser at it. One
  source of truth, so a dual-boot laptop, a phone (via Termux), and a
  desktop all see the same streak, because there is only one streak,
  held by the server.
\end{itemize}

The CLI\textquotesingle s own commands (\texttt{ascend},
\texttt{ascend done ...}, \texttt{ascend levels}, ...) are
\textbf{identical in both modes}. Only
\path|ascend config server <url>| +
\texttt{ascend login} differ. This matters for onboarding: a new user
never needs to think about servers to get a working streak tracker
tonight; they opt into the server later, with no relearning.

\hypertarget{spec-v3-public-release-3-multi-user--auth-server-mode-only}{%
\subsubsection{3. Multi-user \& auth (server mode
only)}\label{spec-v3-public-release-3-multi-user--auth-server-mode-only}}

\begin{verbatim}
CREATE TABLE IF NOT EXISTS users (
  id INTEGER PRIMARY KEY, username TEXT UNIQUE NOT NULL,
  password_hash TEXT NOT NULL,         -- PBKDF2/scrypt, stdlib hashlib, never plaintext, never logged
  timezone TEXT NOT NULL DEFAULT 'UTC',-- IANA name, e.g. 'Asia/Calcutta'. REQUIRED once >1 user exists.
  created_at TEXT NOT NULL
);
CREATE TABLE IF NOT EXISTS api_tokens (
  token TEXT PRIMARY KEY, user_id INTEGER NOT NULL, created_at TEXT NOT NULL, label TEXT
);
\end{verbatim}

\begin{itemize}
\tightlist
\item
  Every existing table (\texttt{entries}, \texttt{node\_state},
  \texttt{missions}, \texttt{packs}, \texttt{parked}, \texttt{settings})
  gains \texttt{user\_id}, and every query is scoped to the
  caller\textquotesingle s user. One SQLite file, many isolated users
  --- same approach Praxis already uses.
\item
  \textbf{Fixed gap:} v1 §6\textquotesingle s \texttt{logical\_date()}
  (the 4 AM rollover) assumed one person\textquotesingle s local time,
  which breaks the moment a server has users in different timezones ---
  "local" becomes meaningless on a server process. Every call to
  \texttt{logical\_date(now)} in server mode becomes
  \path|logical_date(now, user.timezone)|, reading
  \texttt{users.timezone} above. In standalone mode there is exactly one
  implicit user, so this has no effect on you today; it only matters
  once someone else\textquotesingle s account exists.
  \texttt{AUTH\_MODE=single} can default this to the
  system\textquotesingle s local timezone without asking.
\item
  \path|ascend login <server-url>| prompts
  for username/password, stores a long-lived token in
  \path|~/.config/ascend/token| (0600 permissions),
  not the password.
\item
  \texttt{AUTH\_MODE=single} (default for a private deployment, e.g.
  just you on your own Pi) skips login entirely --- one implicit user,
  same as v1/v2. \texttt{AUTH\_MODE=multi} turns on real accounts, for
  anyone who wants to run this for a family, a study group, or the
  public internet. This single switch is what makes the same codebase
  fit "just me" and "a public instance" without two code paths.
\item
  Server binds to the LAN by default (\texttt{-\/-host 0.0.0.0} is
  opt-in, documented with a warning: put it behind a reverse proxy with
  TLS, e.g. Caddy, if exposing it past your LAN).
\end{itemize}

\hypertarget{spec-v3-public-release-4-packs-personal-data-vs-the-public-library}{%
\subsubsection{4. Packs: personal data vs. the public
library}\label{spec-v3-public-release-4-packs-personal-data-vs-the-public-library}}

This is the piece that makes the repo genuinely reusable rather than "my
project with your name filed off":

\begin{verbatim}
ascend/                     <- the installable package (engine, storage, cli, server). No personal data. This is what you publish.
packs/community/            <- curated example packs anyone can use, reviewed, in the repo, e.g.:
  home-hack-lab.toml           networking-101.toml
  dsa-fundamentals.toml        webdev-fullstack.toml
  gamedev-godot.toml           os-internals.toml
~/.config/ascend/packs/     <- a USER's own private packs. Never committed, never shipped. gitignored by default
                                if the user clones the repo to hack on it; irrelevant for a pip-installed user.
\end{verbatim}

\begin{itemize}
\tightlist
\item
  \path|ascend pack install <id>| can
  install from \texttt{packs/community/} (bundled or fetched from the
  repo\textquotesingle s raw URL) or from a local file path. Either way
  it copies into the user\textquotesingle s own data dir; editing your
  copy never touches the shared one.
\item
  Your current cyber pack and personal \texttt{{[}vars{]}} (your
  Pi\textquotesingle s IP, your folder paths) move to
  \path|packs/community/home-hack-lab.toml| as a \emph{generic
  example} (placeholder IP, \texttt{\{hack\_dir\}} already parameterised
  --- it already is, good) --- genuinely useful to a stranger, not wired
  to you specifically. Anything truly yours (your actual IP, your actual
  Telegram token) lives only in your local \texttt{{[}vars{]}} override
  or \texttt{.env}, never in a file you\textquotesingle d commit.
\item
  \texttt{CONTRIBUTING.md} explains how someone submits a new community
  pack via PR: run \texttt{ascend pack lint} and
  \path|ascend pack check-links|, fill in the "done\_when produces a
  real artifact" requirement, no AI-hallucinated URLs (same discipline
  you asked for, now a contribution rule, not just a personal one).
\end{itemize}

\hypertarget{spec-v3-public-release-5-verifiers-as-a-real-plugin-system-this-is-your-customizability-the-triliumcherrytree-equivalent-of-pluginsscripts}{%
\subsubsection{5. Verifiers as a real plugin system (this is your
"customizability," the Trilium/CherryTree equivalent of
plugins/scripts)}\label{spec-v3-public-release-5-verifiers-as-a-real-plugin-system-this-is-your-customizability-the-triliumcherrytree-equivalent-of-pluginsscripts}}

\begin{itemize}
\tightlist
\item
  Built-in verifiers (\texttt{manual}, \texttt{file}, \texttt{command},
  later \texttt{http}/\texttt{git}) ship in \texttt{ascend/verifiers/}.
\item
  \textbf{Third-party verifiers}: any Python file exposing
  \path|verify(spec, vars) -> VerifyResult| (v2
  §7\textquotesingle s interface) dropped into
  \path|~/.config/ascend/verifiers/| is loadable via
  \path|verify = { type = "plugin:my_checker", ... }|. This
  is how someone adds a Codeforces checker, a Duolingo-API checker, a
  smart-home sensor checker --- whatever their own "proficiency" needs
  --- without forking the core engine. Exactly the extension point
  Trilium\textquotesingle s scripting API and
  CherryTree\textquotesingle s plugins play, sized for this project.
\item
  \textbf{Trust is explicit and unchanged in spirit from v2 §7}: a
  plugin is code, so \texttt{ascend pack install} lists any
  \texttt{plugin:} verifiers a pack references, and nothing executes
  until \path|ascend plugin trust <name>|
  --- same reasoning as the \texttt{command} verifier trust step,
  extended to Python plugins, since they run with the same permissions
  you have.
\item
  Document the interface in \path|docs/writing-a-verifier.md| with the
  two built-in examples (\texttt{file}, \texttt{command}) as the
  template, so a contributor can write one without reading the whole
  engine.
\end{itemize}

\hypertarget{spec-v3-public-release-6-http-api-server-mode}{%
\subsubsection{6. HTTP API (server mode)}\label{spec-v3-public-release-6-http-api-server-mode}}

\begin{verbatim}
GET  /api/health                              -> 200 OK, no auth (used by docker/Dockerfile's HEALTHCHECK)
POST /api/login                               -> token
GET  /api/me                                  -> streak, active mission, today status  (mirrors `ascend`)
GET  /api/levels?pack=<id>
POST /api/entries        {learned, next_step, node_id?, proof?}   -> mirrors `ascend done`
POST /api/nodes/{id}/mark {status}
POST /api/packs           (multipart TOML upload) -> lint result
GET  /api/packs, GET /api/packs/{id}
GET  /api/stats, GET /api/heatmap
\end{verbatim}

\begin{itemize}
\tightlist
\item
  Same validation rules as the CLI (20-char minimum, secrets guard, one
  active mission) --- the CLI\textquotesingle s local mode and the
  server must enforce identical rules, so behaviour never depends on
  which mode a user picked. In practice this means \texttt{engine/}
  functions take plain data in and out and neither \texttt{cli/} nor
  \texttt{server/} duplicates logic; they are both thin wrappers around
  the same calls.
\item
  A minimal read-only
  \path|GET /api/public/<username>/streak|
  (opt-in per user, off by default) is the hook for a future "show my
  streak on my GitHub profile" badge --- a nice public-repo touch, not
  required for v1.
\end{itemize}

\hypertarget{spec-v3-public-release-7-nudging-simplified-by-server-mode-supersedes-v2-11bs-sync-workaround}{%
\subsubsection{7. Nudging, simplified by server mode (supersedes v2
§11b\textquotesingle s sync
workaround)}\label{spec-v3-public-release-7-nudging-simplified-by-server-mode-supersedes-v2-11bs-sync-workaround}}

\begin{itemize}
\tightlist
\item
  \textbf{If you run a server:} the server itself is the always-on
  machine (your Pi). It owns the single \texttt{nudge\_time} per user
  and sends the Telegram (or, later, ntfy.sh / Pushover / email) message
  directly --- there is no "sync the DB to the Pi" step anymore, because
  the Pi already holds the only copy. This directly fixes the dual-boot
  problem: neither OS needs to be awake at nudge time, since the server
  isn\textquotesingle t on the laptop at all.
\item
  \textbf{If you don\textquotesingle t run a server (standalone mode):}
  the local desktop-notification timer from v2 §11b still applies, as
  the lightweight fallback for a single machine.
\item
  Nudge channels are pluggable with the same shape as verifiers
  (\path|nudge/channels/telegram.py|, \texttt{ntfy.py},
  \texttt{desktop.py}), so someone self-hosting for a study group could
  add Discord or Matrix without touching the core.
\end{itemize}

\hypertarget{spec-v3-public-release-7b-schema-migrations-across-phases-fixed-gap}{%
\subsubsection{7b. Schema migrations across phases (fixed
gap)}\label{spec-v3-public-release-7b-schema-migrations-across-phases-fixed-gap}}

The schema genuinely changes shape across v1 -\textgreater{} v2
-\textgreater{} v3 (new columns, new tables). Prose saying "add this
column" is not a safe upgrade path for someone\textquotesingle s real
database. Adopt the same pattern Praxis already uses successfully:

\begin{verbatim}
storage/migrations/
  0001_init.sql          -- v1: entries, missions, parked, settings
  0002_packs.sql          -- v2: packs, node_state, entries.node_id/source/complete
  0003_multiuser.sql      -- v3: users, api_tokens, user_id columns, users.timezone
\end{verbatim}

\begin{verbatim}
CREATE TABLE IF NOT EXISTS schema_version (version INTEGER NOT NULL);
\end{verbatim}

On every startup (\texttt{ascend}, \texttt{ascend-server}): read the
current version (0 if the table is empty), apply every migration file
numbered above it, in order, inside one transaction per file, and record
the new version. Idempotent, so running it twice is harmless, and it is
what lets you upgrade from "just the Phase 1 CLI I\textquotesingle m
using today" to the server later without losing or corrupting your
existing streak.

\hypertarget{spec-v3-public-release-7c-sqlite-concurrency-fixed-gap}{%
\subsubsection{7c. SQLite concurrency (fixed
gap)}\label{spec-v3-public-release-7c-sqlite-concurrency-fixed-gap}}

Once a server has more than one writer at a time --- your CLI and the
web PWA both logging around the same moment --- SQLite\textquotesingle s
default mode can block or, in rarer cases, surprise you. Two concrete,
cheap mitigations, not a rewrite:

\begin{itemize}
\tightlist
\item
  \path|PRAGMA journal_mode=WAL;| on every connection at startup ---
  readers no longer block writers.
\item
  The server process is the only thing that opens the DB in server mode
  (the CLI talks to it over HTTP, never opens the file directly once a
  server is configured) --- a single-writer-per-process design, which is
  sufficient at the scale this project is designed for (one person, a
  study group, not internet-scale traffic). If that ever stops being
  true, that is the point to move off SQLite, not before.
\end{itemize}

\hypertarget{spec-v3-public-release-8-deployment}{%
\subsubsection{8. Deployment}\label{spec-v3-public-release-8-deployment}}

\textbf{On your Pi (recommended for you --- always-on, matches your
hardware):}

\begin{verbatim}
# docker-compose.yml
services:
  ascend:
    build: .
    ports: ["8420:8420"]
    volumes: ["./data:/data"]       # SQLite file lives here, back this up
    environment:
      - AUTH_MODE=single            # set to "multi" if others will use this instance
      - TELEGRAM_BOT_TOKEN=${TELEGRAM_BOT_TOKEN}   # from a local .env, gitignored
      - TELEGRAM_CHAT_ID=${TELEGRAM_CHAT_ID}
\end{verbatim}

\texttt{docker compose up -d} on the Pi (arm64 base image, same
pattern Praxis already documents for itself). CLI clients on any machine
set
\path|ascend config server http://<pi-ip>:8420|.

\textbf{For anyone else cloning the repo:} same Dockerfile/compose runs
on an x86 VPS, a Raspberry Pi, or just \texttt{localhost} for testing
--- nothing Pi-specific is hardcoded. Document all three in the
README\textquotesingle s "Deploy" section. A plain \texttt{pip install}
with no server is still the zero-friction path for someone who just
wants the CLI on one machine.

\hypertarget{spec-v3-public-release-9-open-source-polish-checklist}{%
\subsubsection{9. Open-source polish
checklist}\label{spec-v3-public-release-9-open-source-polish-checklist}}

Before making the repo public:

\begin{itemize}
\tightlist
\item
  \texttt{LICENSE} --- MIT is the simplest permissive default and
  matches most CLI tools; CherryTree is GPLv3 and Trilium is AGPLv3,
  which are fine alternatives if you specifically want
  modifications/hosted forks to stay open. Pick one deliberately rather
  than leaving it unset (an unlicensed public repo legally defaults to
  "no rights granted").
\item
  \texttt{README.md} --- what it is (one paragraph), a GIF/screenshot of
  \texttt{ascend} and \texttt{ascend levels}, quickstart
  (\texttt{pip install}, \path|ascend pack install home-hack-lab|,
  \texttt{ascend}), then "Server mode" and "Writing your own pack" as
  later sections. Lead with standalone mode; server mode is "when
  you\textquotesingle re ready," not a prerequisite.
\item
  \texttt{CONTRIBUTING.md} --- how to submit a community pack (section
  4) and a verifier (section 5), code style, how to run tests.
\item
  \texttt{docs/} --- pack format reference (v2 §3 content, cleaned up),
  verifier-writing guide, API reference (section 6), deployment guide
  (section 8).
\item
  \texttt{examples/} --- 2-3 finished example packs across different
  domains (cyber, DSA, a language, a creative skill) so a visitor
  immediately sees this isn\textquotesingle t a cyber-only or CS-only
  tool.
\item
  Tests + CI (GitHub Actions running \texttt{pytest} +
  \texttt{ascend pack lint} on every community pack) --- keeps
  community contributions honest without you manually reviewing every
  TOML line by line.
\item
  A clear "Status" note if you publish before Phase 2 (server) is built,
  so early users know standalone mode is solid and server mode is in
  progress --- honest scoping beats an overpromised README.
\end{itemize}

\hypertarget{spec-v3-public-release-10-build-order-unchanged-principle-dont-let-building-this-become-its-own-overwhelm-spiral}{%
\subsubsection{10. Build order (unchanged principle:
don\textquotesingle t let building this become its own overwhelm
spiral)}\label{spec-v3-public-release-10-build-order-unchanged-principle-dont-let-building-this-become-its-own-overwhelm-spiral}}

\begin{enumerate}
\def\labelenumi{\arabic{enumi}.}
\tightlist
\item
  \textbf{Phase 1 (v2):} standalone engine + CLI + packs +
  levels/milestones. Ship this alone first --- it is already a complete,
  useful, installable tool with zero dependencies.
\item
  \textbf{Phase 1b (v2):} local nudging (desktop notification).
\item
  \textbf{Phase 2 (this file, REVISED): server + a real web client
  together, not server-alone.} A headless API with no UI is not useful
  day to day, and checking progress from a phone matters as much as the
  server itself. Build:

  \begin{itemize}
  \tightlist
  \item
    \texttt{server/} (FastAPI) behind the identical CLI commands, auth,
    the HTTP API (§6), Docker Compose (§8).
  \item
    \texttt{web/}: a small PWA (one static page + a manifest + a service
    worker) hitting the same API: shows the next step, the quick win,
    levels/streak/heatmap, and a form to log \texttt{ascend done}
    equivalents --- a FULL client, not read-only, so the phone alone is
    enough on a day you never open a terminal. Add to Android home
    screen via the browser\textquotesingle s "Install app" prompt
    (standard PWA behaviour, no app-store listing needed).
  \item
    Notification channel (§12 below), default \textbf{ntfy}, since it is
    the least infrastructure for a real Android push.
  \end{itemize}
\item
  \textbf{Phase 3:} plugin verifiers + trust step, community pack
  registry + lint CI, docs/CONTRIBUTING polish, then publish.
\item
  \textbf{Phase 4:} Telegram channel + inline-button quick actions, true
  Web Push from the PWA (no separate app at all), public-facing extras
  (streak badge, etc.).
\end{enumerate}

\hypertarget{spec-v3-public-release-12-interacting-beyond-the-cli-and-real-android-notifications}{%
\subsubsection{12. Interacting beyond the CLI, and real Android
notifications}\label{spec-v3-public-release-12-interacting-beyond-the-cli-and-real-android-notifications}}

\textbf{Two clients, one source of truth, pick per moment, not
permanently.} CLI is fastest when you\textquotesingle re already at a
terminal mid-step (SSH\textquotesingle d into the Pi, running nmap, then
straight to \texttt{ascend done}). The web PWA is the one you open from
your phone to check the streak, see which level you\textquotesingle re
on, or log a step on a day with no terminal in sight. Neither is "the
real one"; both call the identical API from §6, so progress never forks
between them.

\textbf{Why ntfy is the default notification channel, not Telegram and
not Web Push, to start:}

\begin{longtable}[]{@{}P{.22\textwidth}P{.18\textwidth}P{.32\textwidth}P{.18\textwidth}@{}}
\toprule\noalign{}
Option & Android push? & Setup effort & Notes \\
\midrule\noalign{}
\endhead
\bottomrule\noalign{}
\endlastfoot
CLI alone & No & --- & Cannot push anything; only reacts when run \\
\textbf{ntfy} & Yes, real push & Lowest: install the app, subscribe to
one topic, server does a plain HTTP POST & Free hosted \texttt{ntfy.sh}
or self-host later; no account/bot needed \\
Telegram & Yes, real push & Low-medium: create a bot, get a chat id &
Bonus: inline buttons can log a step directly from the notification \\
Web Push (PWA) & Yes, on Android Chrome & Highest: needs HTTPS + VAPID
keys + a reachable server & No separate app installed at all --- the
most "one unified app" feel, worth it once the project is mature \\
\end{longtable}

Ship ntfy first (§7\textquotesingle s pluggable channel design already
anticipated this: \path|nudge/channels/ntfy.py|), add Telegram as an
alternative users can switch to
(\path!ascend config nudge-channel ntfy|telegram|desktop!),
and leave Web Push as a Phase 4 upgrade once the server has a real
domain and TLS anyway.

Do not start Phase 2 until Phase 1 genuinely works for you day to day. A
public repo with an unused server is worse than a private tool with a
real streak --- the whole point of this project is proof that something
got finished.

\hypertarget{spec-v3-public-release-11-prompt-for-the-coding-agent-replaces-v2-13-going-forward}{%
\subsubsection{11. Prompt for the coding agent (replaces v2 §13 going
forward)}\label{spec-v3-public-release-11-prompt-for-the-coding-agent-replaces-v2-13-going-forward}}

\begin{quote}
Read \path|hackstreak-spec-v2.md| for the
engine/pack/level/milestone/streak rules, and
\path|spec-v3-public-release.md| for scope and architecture (v3
supersedes v2 on anything about naming, personal data, or deployment).
Build \textbf{Phase 1 + Phase 1b only, in standalone mode}, as a clean
\texttt{engine/} + \texttt{storage/} + \texttt{cli/} package (already
separated so Phase 2\textquotesingle s server can wrap the same engine
later without a rewrite). Do not build \texttt{server/} yet. Rename the
CLI binary to \texttt{ascend} (or ask me if I\textquotesingle ve chosen
a different name by then). Keep the sample community pack at
\path|packs/community/home-hack-lab.toml| with placeholder
\texttt{{[}vars{]}}, not my real IP or paths. Write tests first. No
dependencies beyond the standard library. Ask me before starting
anything in section 10\textquotesingle s Phase 2 or later.
\end{quote}

\hypertarget{spec-v3-public-release-13-what-verified-means-here-honestly}{%
\subsubsection{13. What "verified" means here,
honestly}\label{spec-v3-public-release-13-what-verified-means-here-honestly}}

This spec is a reasoned design, produced by thinking through the problem
and reviewing it for gaps (§7b\textquotesingle s migration runner,
§7c\textquotesingle s SQLite concurrency note, and §3\textquotesingle s
per-user timezone fix all came from that review pass --- not from
building something and discovering the bug the hard way). It is
\textbf{not verified} in the sense that matters most: no line of it has
run yet, and a design review on paper cannot find everything a real
implementation will surface. That is normal for software, not a flaw
unique to this plan --- treat every "ask me before deviating"
instruction to the coding agent as the mechanism for surfacing the next
real gap, not as a formality.

Two gaps are left deliberately open rather than papered over:

\begin{itemize}
\tightlist
\item
  \textbf{Exact API request/response contracts} (§6) --- only a route
  list exists, not pinned JSON schemas or error codes.
\item
  \textbf{Plugin sandboxing detail} (§5) --- a trust step is specified,
  but the actual sandboxing mechanics are not designed.
\end{itemize}

Both are real LLD work, and writing it now, before Phase 1 (which needs
neither) even exists, would be specifying against a moving target. Do
that work when Phase 2/3 actually starts, informed by how Phase 1
behaved in practice.

\textbf{The actual risk reduction isn\textquotesingle t a longer spec,
it\textquotesingle s this order:} ship the small Phase 1 CLI, write its
tests first and make them pass for real (not just plan to), use it daily
until it earns trust, THEN build the server on a foundation that has
already proven itself. A design that looks complete on paper and a
system that has survived a week of real use are not the same kind of
confidence, and only the second one actually answers "will I fail after
building this."

\hypertarget{spec-v3-public-release-14-feature-audit-what-it-has-what-it-should-have-what-it-deliberately-doesnt}{%
\subsubsection{14. Feature audit: what it has, what it should have, what
it deliberately
doesn\textquotesingle t}\label{spec-v3-public-release-14-feature-audit-what-it-has-what-it-should-have-what-it-deliberately-doesnt}}

\textbf{"Customizable" was always meant on one specific axis: the
roadmap, not the UI.} That axis is genuinely strong:

\begin{longtable}[]{@{}P{.36\textwidth}P{.59\textwidth}@{}}
\toprule\noalign{}
Already specified & Where \\
\midrule\noalign{}
\endhead
\bottomrule\noalign{}
\endlastfoot
Roadmap is editable data (TOML), not code & v2 §3 \\
Levels, gated main path + always-open side lane & v2 §3, §5 \\
Milestone-gated proficiency, not just checkbox completion & v2 §8b \\
Pluggable verifiers (built-in + third-party Python plugins) & v2 §7, v3
§5 \\
Two pack kinds (project vs. proficiency), both AI-draftable or
hand-written & v2 §8b, §9 \\
Numbered-series generator (\texttt{\{from, to\}}) for any
wargame/course with levels & v2 §3 \\
User \texttt{{[}vars{]}} substitution, so one pack works across
different setups & v2 §3 \\
Multi-user, each isolated, one codebase for solo and shared use & v3
§3 \\
Pluggable nudge channels (ntfy/Telegram/desktop, Discord/Matrix addable)
& v3 §7 \\
Same API for CLI and web, so neither is a second-class client & v3 §6,
§12 \\
\end{longtable}

Measured honestly against Trilium/CherryTree, this spec is strong on
\textbf{content/roadmap} extensibility and weak on \textbf{interface}
extensibility and a few \textbf{data-hygiene} basics those tools treat
as table stakes. Below is the honest gap list, split into
what\textquotesingle s worth fixing now (cheap, in the spec, no new
phase) versus what\textquotesingle s correctly deferred.

\hypertarget{spec-v3-public-release-14a-adding-now-cheap-to-spec-real-gaps-no-phase-change}{%
\paragraph{14a. Adding now (cheap to spec, real gaps, no phase
change)}\label{spec-v3-public-release-14a-adding-now-cheap-to-spec-real-gaps-no-phase-change}}

\begin{itemize}
\tightlist
\item
  \textbf{Entries are append-only with no fix-up path.} If the owner
  logs the wrong node or mistypes \texttt{learned}, there was no way to
  correct it. Added:
  \path|ascend edit <entry-id> {[}---learned ...{]} {[}---next-step ...{]}|
  and \texttt{ascend undo} (removes the most recent entry made in the
  last 10 minutes only, to stop it being an easy way to quietly rewrite
  history). Both write an audit row rather than silently overwriting, so
  the streak computation stays honest and inspectable.
\item
  \textbf{\texttt{minutes} was captured but never surfaced.}
  v1\textquotesingle s schema has a \texttt{minutes} column on
  \texttt{entries}; nothing in the CLI ever lets the owner set it or see
  it back. Added: \path|ascend done ... ---minutes N| and
  \texttt{minutes} totals per level/week in \texttt{ascend stats}, so
  "proficiency" has a real time-investment number behind it, not just a
  day count.
\item
  \textbf{No search over months of accumulated notes.} Every
  \texttt{learned} line is a small, real note; after a few months
  that\textquotesingle s a genuine knowledge log, and there was no way
  to find anything in it. Added:
  \path|ascend search "<text>"| (substring
  over \texttt{learned}/\texttt{next\_step}/\texttt{proof}, SQLite
  \texttt{LIKE} is sufficient at this scale --- no search engine needed)
  and the same as \texttt{GET /api/search?q=} on the server.
\item
  \textbf{No visual skill map, despite the whole design being modeled on
  one.} \texttt{hack levels} was specified as a text progress bar; the
  actual Duolingo signature feature --- a path of nodes you can see, not
  just a number --- was never designed. Added to the web PWA
  (§10\textquotesingle s Phase 2 scope, not a new phase): a simple
  level-path view, one row per level, nodes as circles
  (locked/open/done/milestone), reusing data already returned by
  \texttt{GET /api/levels}. No separate visual-design pass is scheduled
  for this now --- ship the functional version in Phase 2; polishing it
  is Phase 4.
\item
  \textbf{No theming, despite customization being the explicit goal.}
  Trilium and CherryTree both ship themes; this spec\textquotesingle s
  web PWA had none specified. Added: CSS custom properties for color
  tokens, a light/dark toggle defaulting to
  \texttt{prefers-color-scheme}, stored in \texttt{localStorage} (view
  preference only, never synced --- losing it on a new device is an
  acceptable trade-off for the amount of code it saves). This is
  near-zero marginal cost once the PWA exists, so it\textquotesingle s
  folded into Phase 2, not deferred.
\item
  \textbf{One global nudge time couldn\textquotesingle t fit a varied
  schedule.} The owner\textquotesingle s week already has lab-night vs.
  free-half-day rhythm (see the original hack-lab planning); one fixed
  \texttt{nudge\_time} doesn\textquotesingle t reflect that. Added: an
  optional \texttt{nudge\_time} override per mission
  (\path|ascend mission start <pack> ---nudge-time 21:30|),
  falling back to the global setting when unset.
\end{itemize}

\hypertarget{spec-v3-public-release-14b-real-gaps-deliberately-not-added-with-the-actual-reason-not-just-later}{%
\paragraph{14b. Real gaps, deliberately NOT added (with the actual
reason, not just
"later")}\label{spec-v3-public-release-14b-real-gaps-deliberately-not-added-with-the-actual-reason-not-just-later}}

\begin{itemize}
\tightlist
\item
  \textbf{Full theming/branding system, i18n, accessibility audit.}
  These matter a lot for a public tool with many users, and genuinely
  little for the first real user (you, in English, on your own devices).
  Building them now would be designing for an audience that
  doesn\textquotesingle t exist yet. Revisit once outside contributors
  actually show up --- their needs should shape this, not a guess made
  in advance.
\item
  \textbf{Public-signup hardening (rate limiting, invite-only mode,
  email verification).} \texttt{AUTH\_MODE=multi} as specified is
  adequate for a study group you personally invite; it is NOT specified
  to be safe as an open-signup public service. Don\textquotesingle t
  flip a server to open multi-user on the public internet without a real
  security pass first --- this is a correctness gap I\textquotesingle m
  flagging, not one I\textquotesingle m fixing by writing more spec
  prose at this stage, because the right fix depends on
  who\textquotesingle s actually using it.
\item
  \textbf{Offline-first PWA (background sync queue, conflict resolution
  for offline edits).} Real engineering, real value eventually, but the
  server is reachable from your phone over the same network as the Pi
  almost all the time; building a sync-conflict system for an offline
  scenario that rarely occurs yet would be solving a problem you
  don\textquotesingle t have.
\item
  \textbf{Social/leaderboard features (compare streaks with friends,
  encouragement, comments).} Explicitly against the original design
  goal: the whole point was escaping overwhelm and open-ended systems,
  not adding a second source of social pressure. If you want this later
  for a study group, it\textquotesingle s a small addition on top of the
  public-streak hook already in §6 --- but it\textquotesingle s opt-in
  territory, not a default.
\item
  \textbf{Account export/deletion (GDPR-style "give me all my data" /
  "delete everything").} Only matters once real
  strangers\textquotesingle{} data is on a server you run. Required
  before genuinely public multi-user hosting, irrelevant before that.
  Flagged here so it isn\textquotesingle t forgotten, not built now.
\end{itemize}

\textbf{Where this leaves you:} the roadmap/content side --- the part
you specifically asked not to be hardcoded --- is the best-covered part
of the design. The interface side had real, fixable gaps (no edit/undo,
no search, no visual map, no theming), now folded into the existing
Phase 2 scope at no cost to the build order. The remaining gaps are real
but correctly out of scope until there\textquotesingle s an actual
audience whose needs would shape them better than a guess would.
